The rank of~$M$ can be viewed as a combinatorial property, in that matrix multiplication is not commutative, and even the rank of a matrix product can depend on the order of the multiplied matrices.
Extending the results of \autoref{sec-toolbox} but aiming at a more structural result, we address a more general question in this section: is there a specific vector space, perhaps a joint invariant subspace of the generating matrices, whose dimension tells us something about the rank of $M$? 

In total DFAs, every row $[q]^T$ is maximal, so $\spn{\MRow} = \mathbb{R}^Q$. In \cite[Criterion~1]{DBLP:journals/isci/BerlinkovS16}, the following result was proved. A different proof was independently and concurrently provided in \cite[\S8]{Voynov2015Compact}, see also \cite{Protasov2021}\footnote{ \cite[Theorem 2]{Protasov2021} states that if the rank of a total DFA is greater than one, then there exists a non-trivial joint invariant subspace of its generating matrices. However, such a subspace exists for all total DFAs regardless of their rank~\cite[Corollary 4]{Protasov2017}. 
Hence, to characterise total DFAs of rank greater than one, we have to consider the existence of some specific joint invariant linear subspace, for example the one in the statement of \autoref{thm:alg-criterion-dfas}.}.

\begin{thm}[Algebraic synchronization criterion for total DFAs]\label{thm:alg-criterion-dfas}
    If $M$ is a total DFA, we have 
    $\spn{\alpha^T M(w) \mid w \in \Sigma^*} = \mathbb{R}^Q$
    if and only if $r=1$.
\end{thm}

% Fix indentation
\noindent
It is thus reasonable to ask if this statement can be generalised to the case where $r > 1$ and $M$ is an arbitrary unambiguous monoid morphism. The theorem below provides such a generalisation, using the results obtained above. In particular, it implies that for total DFAs $\dim \spn{\alpha^T M(w) \mid w \in \Sigma^*} \le n - r + 1$. To be consistent with the previous results, we formulate the statements for columns, but an analogous symmetric version for rows, involving $\spn{\alpha^T M(w) \mid w \in \Sigma^*}$ and $\spn{\MRow}$ instead of, respectively, $V$ and $\spn{\MCol}$ also holds true. 

\begin{thm}\label{thm:algebraic}
Define $V := \spn{M(w) \beta \mid w \in \Sigma^*}$.
We have:
\begin{enumerate}
\item[(a)] $V \subseteq \spn{\MCol}$.
\item[(b)] $\dim V + r-1 \le \dim \spn{\MCol}$.
\item[(c)] $V = \spn{\MCol}$ if and only if $r=1$.
\end{enumerate}
\end{thm}

\begin{proof}
Towards item~(a), it suffices to show that $\spn{\MCol}^\bot \subseteq V^\bot$.
Let $x^\top \in \spn{\MCol}^\bot$.
For every word $u \in \Sigma^*$ of minimum rank, by \autoref{thm:Cesari}, each column of~$M(u)$ is in $\spn{\MCol}$.
Thus, $x^\top M(u) \beta = 0$ holds for all $u \in \Sigma^*$ of minimum rank.
By \autoref{prop:compact}, it follows that $x^\top M(w) \beta = 0$ holds for all $w \in \Sigma^*$, i.e., $x^\top \in V^\bot$.

Towards item~(b), let $u \in \Sigma^*$ be of minimum rank.
By \autoref{thm:Cesari}, there are $r$ maximal columns $[C_1], \ldots, [C_r]$ and $r$ maximal rows $[R_1]^T, \ldots, [R_r]^T$ such that $M(u) = \sum_{i=1}^r [C_i] [R_i]^T$.
We conclude from item~(a) that $\spn{V, [C_1], \ldots, [C_r]} \subseteq \spn{\MCol}$.
Hence, it suffices to show that $\dim V + r-1 = \dim \spn{V, [C_1], \ldots, [C_{r-1}]}$.

We have
\[
M(u u) \ = \ \left( \sum_{i=1}^r [C_i] [R_i]^T \right) \left( \sum_{j=1}^r [C_j] [R_j]^T \right)
\ = \ \sum_{1 \le i,j \le r} [C_i] [R_i]^T [C_j] [R_j]^T\,.
\]
By unambiguousness, $[R_i]^T [C_j] \le 1$ for all $1 \le i,j \le r$.
For each~$i$ there is at most one~$j$ with $[R_i]^T [C_j] = 1$ as otherwise $[R_j]^T + [R_{j'}]^T$ for some $j' \ne j$ would appear in~$M(u u)$, contradicting the maximality of~$[R_j]^T$.
On the other hand, for each~$i$ there is at least one~$j$ with $[R_i]^T [C_j] = 1$ as otherwise $M(u u)$ has only $[C_{i'}]$ with $i' \ne i$ as nonzero columns, contradicting the fact that $M(u u)$ has rank~$r$.
Thus, for each~$i$ there is exactly one~$j$ with $[R_i]^T [C_j] = 1$.
With a symmetric argument, for each~$j$ there is exactly one~$i$ with $[R_i]^T [C_j] = 1$.
It follows that there is permutation $\pi: \{1, \ldots, r\} \to \{1, \ldots, r\}$ such that $[R_i]^T [C_j] = 1$ if and only if $i = \pi(j)$.

To show that $\dim V + r-1 = \dim \spn{V, [C_1], \ldots, [C_{r-1}]}$, it suffices to show for each $j \in \{1, \ldots, r-1\}$ that $[C_j] \not\in \spn{V, [C_1], \ldots, [C_{j-1}]}$.
To this end, let $1 \le j < r$ and define $x^T := [R_{\pi(j)}]^T - [R_{\pi(r)}]^T$.
Since $[R_{\pi(j)}]^T$ and $[R_{\pi(r)}]^T$ are maximal rows, for all $w \in \Sigma^*$ also $[R_{\pi(j)}]^T M(w)$ and $[R_{\pi(r)}]^T M(w)$ are maximal rows, implying that $x^T M(w) \beta = \mrw - \mrw = 0$; i.e., $x^T$ is in the orthogonal complement of~$V$.
Further, for all $i \in \{1, \ldots, j-1\}$ we have $x^T [C_i] = [R_{\pi(j)}]^T [C_i] - [R_{\pi(r)}]^T [C_i] = 0-0 = 0$; i.e., $x^T$~is orthogonal also to~$[C_i]$.
But $x^T$~is not orthogonal to~$[C_j]$, as $x^T [C_j] = [R_{\pi(j)}]^T [C_j] - [R_{\pi(r)}]^T [C_j] = 1 - 0 = 1$.
This completes the proof of item~(b).

Towards item~(c), if $r>1$ we have $V \subsetneq \spn{\MCol}$ by item~(b).
Let $r=1$.
By item~(a) we have $V \subseteq \spn{\MCol}$.
Towards the opposite inclusion, let $y \in \MCol$.
Since $r=1$, there is $w \in \Sigma^*$ such that $y$ is the only nonzero column of~$M(w)$.
Thus, $y$~is a multiple of $M(w) \beta$.
Hence, $y \in V$.
\end{proof}

\begin{rem}
The inequality in \autoref{thm:algebraic}~(b) and its row version can be both strict if $r>1$, even in total DFAs. In \autoref{ex:snap1}, $V$ is spanned by $\begin{pmatrix} 1 & 1 & 1 & 1\end{pmatrix}^T$, but 
\[\dim \spn{\MCol} = 3 > 1 + 2 - 1 = \dim V + r-1.\]
For the row version, $\spn{\alpha^T M(w) \mid w \in \Sigma^*}$ is spanned by 
$\begin{pmatrix}1 & 0 & 1 & 0\end{pmatrix}$ and
$\begin{pmatrix}0 & 1 & 0 & 1\end{pmatrix}$, but 
\[\dim \spn{\MRow} = 4 > 2 + 2 - 1 = \dim \spn{\alpha^T M(w) \mid w \in \Sigma^*} + r - 1.\]
\end{rem}
